\label{sec:introduction}

Algorithmic redistricting is not merely a theoretical alternative to
political map-drawing.  It is a system that can be operationalized---run
by identified actors, subject to identified oversight procedures, and
producing outputs that enter an identified legal process.  Abstracting
away from implementation, as much of the literature does, leaves open the
critical question of whether algorithmic redistricting is politically and
institutionally feasible.

This paper answers that question by developing concrete case studies for
three adoption pathways.

\textbf{The Commission Pathway} (Case Study 1 --- Arizona AIRC) is the path
of least resistance where an independent redistricting commission already
exists.  The commission could treat the algorithm's output as a starting
plan that satisfies mathematical criteria, then apply its community-of-interest
adjustment authority to refine boundaries before public adoption.

\textbf{The Statute Pathway} (Case Study 2 --- California CRC) involves
legislative or commission action in a state that has already adopted
multi-criterion redistricting criteria.  California's criteria---population
equality, VRA compliance, county contiguity, compactness, and
proportionality---map closely onto the \texttt{redist} system's weights and
constraints.  A California-style adoption would require demonstrating that
each criterion is satisfied algorithmically, a task the paper shows is
tractable.

\textbf{The Court-Order Pathway} (Case Study 3 --- North Carolina
post-\textit{Harper}) is available where state courts have found
constitutional violations in enacted maps.  The court-order pathway places
the algorithm at the center of the remedial process: a special master
appointed by the court runs \texttt{redist}, the output becomes the
remedial map, and the burden shifts to the challenger to demonstrate a
constitutional defect in the output.

These three pathways are not mutually exclusive.  A state might adopt
statutory criteria that require algorithmic verification, or a court might
order a commission to use an algorithm.  But treating them as distinct
ideal types allows us to identify the institutional requirements, oversight
structures, and legal challenges specific to each.

The paper is organized as follows.  Sections~2--4 develop the three case
studies in detail.  Section~5 provides a practical workflow table that
identifies, at each stage of the redistricting process, who runs which
command, when, under what oversight, and with what review opportunity.
Section~6 addresses challenges common to all three pathways: data preparation,
public transparency, and the challenge process.  Section~7 concludes.

Throughout, we use the \texttt{redist} system~\citep{dellaLibera2026recursive}
as the reference implementation.  The commands cited in the workflow table
are drawn from the \texttt{redist} CLI reference~\citep{redistCliQuickstart}.
The legal analysis draws on~\citet{dellaLibera2026legal}.
